    \chapter{Symmetric, Hermitian and Unitary Operators}
    \section{Operators and Transpose}
    \begin{definitionenv}
        Suppose $V$ is a vector space over $\mathcal{K}$. A linear map $T: V \to V$ is called a \textbf{linear operator} on $V$. 
    \end{definitionenv}
    \begin{definitionenv}
    Let $V$ be a finite dimensional vector space over field $\mathcal{K}$. The \textbf{dual space} of $V$ is denoted by $V^*$ and is defined as
    \[
    V^*=\mathcal{L}(V,\mathcal{K}),
    \]
    i.e., the vector space of all linear maps from $V$ to $\mathcal{K}$. An element of $V^*$ is called a \textbf{linear functional} on $V$.
    \end{definitionenv}

    \begin{exampleenv}
    \leavevmode
    \begin{enumerate}
        \item Let $V$ be a finite dimensional vector space over field $\mathcal{K}$ with a scalar product. Suppose $v_0\in V$. Define
        \[
        f_{v_0}:V\to \mathcal{K},\qquad f_{v_0}(v)=\langle v,v_0\rangle
        \quad \text{for any } v\in V.
        \]
        Then $f_{v_0}$ is a linear functional on $V$.

        \item Let $V$ be a finite dimensional vector space over $\mathbb{C}$ with a hermitian product. Suppose $v_0\in V$. Define
        \[
        f_{v_0}:V\to \mathbb{C},\qquad f_{v_0}(v)=\langle v,v_0\rangle
        \quad \text{for any } v\in V.
        \]
        Then $f_{v_0}$ is a linear functional on $V$.

        However, the map
        \[
        g_{v_0}:V\to \mathbb{C},\qquad g_{v_0}(v)=\langle v_0,v\rangle
        \quad \text{for any } v\in V
        \]
        is, in general, not a linear functional on $V$ by the conjugate linearity in the second argument of the hermitian product.
    \end{enumerate}
    \end{exampleenv}

    \begin{lemmaenv}
    Let $V$ be a finite dimensional vector space over field $\mathcal{K}$.
    \begin{enumerate}
        \item $\dim(V)=n$ if and only if $\dim(V^*)=n$. Thus $V$, $V^*$, and $\mathcal{K}^n$ are all isomorphic to each other as vector spaces over $\mathcal{K}$.

        \item Suppose $\{v_1,\dots,v_n\}$ is a basis of $V$. Then there exists a unique basis $\{f_1,\dots,f_n\}$ of $V^*$ such that
        \[
        f_i(v_j)=\delta_{ij}
        \qquad
        \text{for any } i,j\in\{1,\dots,n\},
        \]
        where $\delta_{ij}$ is the Kronecker delta function.

        \item Suppose in addition that $V$ is equipped with a non-degenerate scalar product. Then the map
        \[
        \Phi:V\to V^*,\qquad \Phi(v)=f_v,
        \]
        where
        \[
        f_v(w)=\langle w,v\rangle
        \qquad
        \text{for any } w\in V,
        \]
        is an isomorphism of vector spaces over $\mathcal{K}$.
    \end{enumerate}
    \end{lemmaenv}

    \begin{proofenv}
    \leavevmode
    \begin{enumerate}
        \item Assume $\dim(V)=n$, and let $\{v_1,\dots,v_n\}$ be a basis of $V$. For each $i\in\{1,\dots,n\}$, define
        \[
        f_i\!\left(\sum_{j=1}^n a_jv_j\right)=a_i.
        \]
        This is well-defined because every vector in $V$ has a unique expression in the basis $\{v_1,\dots,v_n\}$, and each $f_i$ is linear.

        We claim that $\{f_1,\dots,f_n\}$ is a basis of $V^*$. First, it is linearly independent: if
        \[
        c_1f_1+\cdots+c_nf_n=0,
        \]
        then evaluating at $v_j$ gives
        \[
        0=(c_1f_1+\cdots+c_nf_n)(v_j)=c_j,
        \]
        so all $c_j=0$.

        Next, it spans $V^*$. Let $f\in V^*$. Set
        \[
        a_i=f(v_i)\qquad (i=1,\dots,n).
        \]
        Then for any $v=\displaystyle\sum_{j=1}^n x_jv_j$,
        \[
        f(v)
        =
        f\!\left(\sum_{j=1}^n x_jv_j\right)
        =
        \sum_{j=1}^n x_jf(v_j)
        =
        \sum_{j=1}^n x_ja_j
        =
        \sum_{j=1}^n a_jf_j(v).
        \]
        Hence
        \[
        f=\sum_{j=1}^n a_jf_j.
        \]
        Therefore $\{f_1,\dots,f_n\}$ is a basis of $V^*$, so $\dim(V^*)=n$.

        Conversely, if $\dim(V^*)=n$, then applying the same result to the vector space $V^*$ gives
        \[
        \dim((V^*)^*)=n.
        \]
        Since $V$ is finite dimensional, the canonical map from $V$ to $(V^*)^*$ is an isomorphism, hence
        \[
        \dim(V)=\dim((V^*)^*)=n.
        \]
        So $\dim(V)=n$ if and only if $\dim(V^*)=n$.

        Since any $n$-dimensional vector space over $\mathcal{K}$ is isomorphic to $\mathcal{K}^n$, it follows that $V\cong V^*\cong \mathcal{K}^n$.

        \item Existence has already been established in part (1): the functionals $f_i$ defined by
        \[
        f_i\!\left(\sum_{j=1}^n a_jv_j\right)=a_i
        \]
        satisfy
        \[
        f_i(v_j)=\delta_{ij}.
        \]

        We now show uniqueness. Suppose $\{g_1,\dots,g_n\}$ is another basis of $V^*$ such that
        \[
        g_i(v_j)=\delta_{ij}
        \qquad
        \text{for all } i,j.
        \]
        Let $v=\displaystyle \sum_{j=1}^n a_jv_j$. Then
        \[
        g_i(v)
        =
        g_i\!\left(\sum_{j=1}^n a_jv_j\right)
        =
        \sum_{j=1}^n a_jg_i(v_j)
        =
        \sum_{j=1}^n a_j\delta_{ij}
        =
        a_i
        =
        f_i(v).
        \]
        Since this holds for every $v\in V$, we have $g_i=f_i$ for all $i$. Hence the basis is unique.

        \item First, $\Phi$ is linear. For any $u,v,w\in V$ and $a,b\in \mathcal{K}$,
        \[
        \Phi(au+bv)=f_{au+bv},
        \]
        and for every $w\in V$,
        \[
        f_{au+bv}(w)
        =
        \langle w,au+bv\rangle
        =
        a\langle w,u\rangle+b\langle w,v\rangle
        =
        af_u(w)+bf_v(w).
        \]
        Hence
        \[
        \Phi(au+bv)=a\Phi(u)+b\Phi(v).
        \]

        Next, $\Phi$ is injective. If $\Phi(v)=0$, then $f_v=0$, i.e.
        \[
        \langle w,v\rangle=0
        \qquad
        \text{for all } w\in V.
        \]
        By non-degeneracy of the scalar product, this implies $v=0$.

        Finally, since $\dim(V)=\dim(V^*)$ by part (1), a linear injective map
        \[
        \Phi:V\to V^*
        \]
        between finite dimensional vector spaces of the same dimension must be surjective. Therefore $\Phi$ is an isomorphism.
    \end{enumerate}
    \end{proofenv}
    % \begin{lemmaenv}
    %     Suppose $V$ is a finite-dimensional vector space over $\mathcal{K}$ with a non-degenerate scalar product $\langle \cdot, \cdot \rangle$. Then the map $v\mapsto f_v$ defined by $f_v(w) = \langle v,w \rangle$ is a vector space isomorphism from $V$ to $V^*$.
    % \end{lemmaenv}
    % \begin{proofenv}
    %     We first show that the map is linear. For any $\alpha, \beta \in \mathcal{K}$ and $v_1, v_2, w \in V$, we need to show $(\alpha v_1 + \beta v_2) \mapsto \alpha f_{v_1}(w) + \beta f_{v_2}(w)$ or namely $f_{\alpha v_1 + \beta v_2}(w) = \alpha f_{v_1}(w) + \beta f_{v_2}(w)$. This is straightforward since
    %     \[
    %     f_{\alpha v_1 + \beta v_2}(w) = \langle \alpha v_1 + \beta v_2, w \rangle = \alpha \langle v_1, w \rangle + \beta \langle v_2, w \rangle = \alpha f_{v_1}(w) + \beta f_{v_2}(w).
    %     \]
    %     We then show that the map is injective. Suppose $f_v = 0$ for some $v \in V$. Then we have $\langle v, w \rangle = 0$ for all $w \in V$. By the non-degeneracy of the scalar product, we have $v=0$.

    %     Since we know $\dim V = \dim V^*$, the map is also surjective and thus a vector space isomorphism.
    % \end{proofenv}
    \begin{theoremenv}\label{transpose existence}
        Suppose $V$ is a finite-dimensional vector space over $\mathcal{K}$ with a non-degenerate scalar product $\langle \cdot, \cdot \rangle$ and $T: V \to V$ is a linear operator. Then there exists a unique linear operator $T^T: V \to V$ such that
        \[
        \langle T(\mathbf{v}), \mathbf{w} \rangle = \langle \mathbf{v}, T^T (\mathbf{w}) \rangle, \quad \forall \mathbf{v}, \mathbf{w} \in V.
        \]
        We call $T^T$ the \textbf{transpose} of $T$.
    \end{theoremenv}
    \begin{proofenv}
        Let $L: V \to \mathcal{K}$ be $L(\mathbf{v}) = \langle T(\mathbf{v}), \mathbf{w} \rangle$ for any $\mathbf{v} \in V$. Note $L$ is a linear map from $V$ to $\mathcal{K}$ and thus a linear functional. Note by the above lemma, we know there exists a unique $\tilde{w} \in V$ such that $L(x)=f_{\tilde{w}}(x) = \langle x, \tilde{w} \rangle$. Let $T^T(\mathbf{w}) = \tilde{w}$.

        Suppose $L_1(\mathbf{v}) = \langle T(\mathbf{v}), \mathbf{w}_1 \rangle$ and $\tilde{w}_1 \in V$ is the unique vector in $V$ such that $L_1 = f_{\tilde{w}_1}$. Suppose $L_2(\mathbf{v}) = \langle T(\mathbf{v}), \mathbf{w}_2 \rangle$ and $\tilde{w}_2 \in V$ is the unique vector in $V$ such that $L_2 = f_{\tilde{w}_2}$.

        Note $L' := \alpha L_1 + \beta L_2$ satisfies
        \[
        L'(\mathbf{v}) = \alpha \langle T(\mathbf{v}), \mathbf{w}_1 \rangle + \beta \langle T(\mathbf{v}), \mathbf{w}_2 \rangle = \langle T(\mathbf{v}), \alpha \mathbf{w}_1 + \beta \mathbf{w}_2 \rangle.
        \]
        Thus $T^T(\alpha \mathbf{w}_1 + \beta \mathbf{w}_2) = \alpha T^T(\mathbf{w}_1) + \beta T^T(\mathbf{w}_2)$ and thus $T^T$ is linear. We still need to verify $\langle T(\mathbf{v}), \mathbf{w} \rangle = \langle \mathbf{v}, T^T (\mathbf{w}) \rangle$ for all $\mathbf{v}, \mathbf{w} \in V$. This is straightforward since $L(\mathbf{v}) = \langle T(\mathbf{v}), \mathbf{w} \rangle = f_{\tilde{w}}(\mathbf{v}) = \langle \tilde{w}, \mathbf{v} \rangle = \langle T^T(\mathbf{w}), \mathbf{v} \rangle$.
    \end{proofenv}
    \begin{remarkenv}
        We often just write $\langle T(\mathbf{v}), \mathbf{w} \rangle = \langle \mathbf{v}, T^T (\mathbf{w}) \rangle$, $\langle T\mathbf{v}, \mathbf{w} \rangle = \langle \mathbf{v}, T^T \mathbf{w} \rangle$ or $\langle A\mathbf{v}, \mathbf{w} \rangle = \langle \mathbf{v}, A^T \mathbf{w} \rangle$ where $A$ is a general operator (instead of a matrix) on $V$.
    \end{remarkenv}
    \begin{definitionenv}
        Suppose $V$ is a finite-dimensional vector space over field $\mathcal{K}$ and $T$ is an operator on $V$. Suppose $\mathcal{B}$ is a basis for $V$ and let $\mathcal{M}$ be the matrix representation of $T$ with respect to basis $\mathcal{B}$, then we define the determinant of operator $T$ to be the determinant of matrix $\mathcal{M}$, i.e., $\det(T) = \det(\mathcal{M})$.
    \end{definitionenv}
    \begin{remarkenv}
        Note that the determinant of an operator is well-defined since the value of the determinant of the matrix representation of $T$ does not depend on the choice of basis $\mathcal{B}$: suppose that if a different basis $\mathcal{B}'$ is chosen and $\mathcal{M}'$ is the matrix representation of $T$ with respect to $\mathcal{B}'$, then we have $\mathcal{M}' = P^{-1} \mathcal{M} P$ for some invertible matrix $P$ and thus $$\det(\mathcal{M}') = \det(P^{-1} \mathcal{M} P) = \det(P^{-1}) \det(\mathcal{M}) \det(P) = \det(\mathcal{M}).$$
    \end{remarkenv}
    \begin{theoremenv}
        Suppose $V,\mathcal{K}, \langle \cdot, \cdot \rangle$ satisfy the condition in Theorem~\ref{transpose existence}. Let $A,B$ be operators on $V$ and $c\in \mathcal{K}$. Then
        \begin{enumerate}
            \item $(A+B)^T = A^T + B^T$.
            \item $(cA)^T = c A^T$.
            \item $(AB)^T = B^T A^T$, where $AB$ means the composite operator: $V\to V: AB(x) = (A\circ B)(x) = A(B(x))$.
            \item $(A^T)^T = A$.
        \end{enumerate}
    \end{theoremenv}
    \section{Symmetric, Hermitian and Unitary Operators}
    \begin{definitionenv}
        Suppose $V,\mathcal{K}, \langle \cdot, \cdot \rangle$ satisfy the condition in Theorem~\ref{transpose existence} and $A$ is an operator on $V$. We say $A$ is \textbf{symmetric} if $A^T = A$.
    \end{definitionenv}
    \begin{lemmaenv}\label{unique vector for functional}
        Suppose $V$ is a finite-dimensional vector space over $\C$ with a positive definite hermitian product $\langle \cdot, \cdot \rangle$. Given a functional $L$ on $V$, there exists a unique vector $\tilde{w} \in V$ such that $L(x) = \langle x, \tilde{w} \rangle$ for all $x\in V$.
    \end{lemmaenv}
    \begin{theoremenv}
        Suppose $V, \C, \langle \cdot, \cdot \rangle$ satisfy the condition in Lemma~\ref{unique vector for functional} and $A$ is an operator on $V$. Then there exists a unique operator $A^*$ on $V$ such that $\langle A x, y \rangle = \langle x, A^* y \rangle$ for all $x,y \in V$. We call $A^*$ the \textbf{adjoint} or \textbf{complex transpose}, or the \textbf{conjugate transpose} of $A$.
    \end{theoremenv}
    \begin{exampleenv}
        Let $V = \C^n$ with the standard hermitian product $\langle x, y \rangle = \displaystyle \sum_{i=1}^n x_i \overline{y_i}$. For any $A \in \M_{n \times n}(\C)$, we have $$\langle A x, y \rangle = x^T A^T \overline{y} = x^T \overline{ \left( \overline{A^T} y \right) } = \langle x, \overline{A^T} y \rangle.$$ Thus $A^* = \overline{A^T}$.
    \end{exampleenv}
    \begin{definitionenv}
        Suppose $V, \C, \langle \cdot, \cdot \rangle$ satisfy the condition in Lemma~\ref{unique vector for functional}. An operator $A$ on $V$ is called \textbf{self-adjoint} or \textbf{hermitian} if $A^* = A$. Particularly, if $A$ is hermitian, then $\langle A x, y \rangle = \langle x, A y \rangle$ for all $x,y \in V$.
    \end{definitionenv}
    \begin{theoremenv}
        Suppose $V, \C, \langle \cdot, \cdot \rangle$ satisfy the condition in Lemma~\ref{unique vector for functional}. Let $A$ and $B$ be operators on $V$ and $\alpha \in \C$. Then
        \begin{enumerate}
            \item $(A+B)^* = A^* + B^*$.
            \item $(\alpha A)^* = \overline{\alpha} A^*$.
            \item $(AB)^* = B^* A^*$, where $AB$ means the composite operator: $V\to V: AB(x) = (A\circ B)(x) = A(B(x))$.
            \item $(A^*)^* = A$.
        \end{enumerate}
    \end{theoremenv}
    \begin{definitionenv}\label{def: unitary operator}
        Suppose $V$ is a finite-dimensional vector space over $\R$ with a positive definite scalar product $\langle \cdot, \cdot \rangle$. An operator $A$ on $V$ is called (real-)\textbf{unitary} or \textbf{orthogonal} if $\langle A x, A y \rangle = \langle x, y \rangle$ for all $x,y \in V$.
    \end{definitionenv}
    \begin{theoremenv}
        Suppose $V$, $\R$, $\langle \cdot, \cdot \rangle$ satisfy the condition in Definition~\ref{def: unitary operator}. Let $A$ be an operator on $V$. Then the following conditions on $A$ are equivalent:
        \begin{enumerate}[label=(\alph*)]
            \item $A$ is real-unitary, i.e., $\langle A x, A y \rangle = \langle x, y \rangle$ for all $x,y \in V$.
            \item $A$ preserves norm of vectores, i.e., $\|A x\| = \|x\|$ for all $x \in V$.
            \item If $v\in V$ is a unit vector, then $A v$ is also a unit vector.
            \item $A^T A = I$, where $I$ is the identity operator on $V$.
        \end{enumerate}
    \end{theoremenv}
    \begin{proofenv}
        (a) $\Rightarrow$ (b): by definition $\| Av \| = \sqrt{\langle Av, Av \rangle} = \sqrt{\langle v, v \rangle} = \|v\|$ for any $v\in V$.

        (b) $\Rightarrow$ (c): trivial.

        (d) $\Rightarrow$ (a): for any $x,y \in V$, we have $\langle A x, A y \rangle = \langle x, A^T A y \rangle = \langle x, y \rangle$.

        (b) $\Rightarrow$ (a): for any $x,y \in V$, we have
        \[
        \begin{aligned}
            \langle A x, A y \rangle &= \frac14 \left( \|A x + A y\|^2 - \|A x - A y\|^2 \right)\\  &= \frac{1}{4} \left( \|x + y\|^2 - \|x - y\|^2 \right) = \langle x, y \rangle.
        \end{aligned}
        \]

        (c) $\Rightarrow$ (b): 
        
        Case 1: if $x=0$, then $\|A x\| = \|0\| = \|x\|$. 
        
        Case 2: if $x\neq 0$, then $\dfrac{x}{\|x\|}$ is a unit vector and thus
        $$ \left\| A\frac{x}{\|x\|}\right\| ^2 = 1 = \langle \frac{x}{\|x\|}, \frac{x}{\|x\|} \rangle = \frac1{\|x\|^2} \langle x, x \rangle = \frac{\|x\|^2}{\|x\|^2} = 1. $$
        Thus $\|A x\| = \|x\|$.

        (a) $\Rightarrow$ (d): for any $x,y \in V$, we have
        $$ \langle Ax, A y \rangle = \langle x,y \rangle = \langle x, A^T A y \rangle \implies \langle x, (A^T A - I) y \rangle = 0. $$
        By the non-degeneracy of the scalar product, we have $(A^T A - I) y = 0$ for all $y\in V$ and thus $A^T A = I$.
    \end{proofenv}
    \begin{exampleenv}
        Let
        $$ R(\theta) = \begin{bmatrix}
            \cos \theta & -\sin \theta \\
            \sin \theta & \cos \theta
        \end{bmatrix}. $$
        For a vector $x\in \R^2$, $R x$ rotated counterclockwise by $\theta$. We also call such a matrix unitary if the associated operator is unitary.
    \end{exampleenv}
    \begin{definitionenv}\label{def: complex unitary operator}
        Suppose $V$ is a finite-dimensional vector space over $\C$ with a positive definite hermitian product $\langle \cdot, \cdot \rangle$. An operator $A$ on $V$ is called (complex-)\textbf{unitary} if $\langle A x, A y \rangle = \langle x, y \rangle$ for all $x,y \in V$.
    \end{definitionenv}
    \begin{theoremenv}
        Suppose $V$, $\C$, $\langle \cdot, \cdot \rangle$ satisfy the condition in Definition~\ref{def: complex unitary operator}. Then the following conditions on an operator $A$ on $V$ are equivalent:
        \begin{enumerate}[label=(\alph*)]
            \item $A$ is complex-unitary, i.e., $\langle A x, A y \rangle = \langle x, y \rangle$ for all $x,y \in V$.
            \item $A$ preserves norm of vectores, i.e., $\|A x\| = \|x\|$ for all $x \in V$.
            \item If $v\in V$ is a unit vector, then $A v$ is also a unit vector.
            \item $A^* A = I$, where $I$ is the identity operator on $V$.
        \end{enumerate}
    \end{theoremenv}
    \begin{propositionenv}
        Let $V$ be a finite-dimensional vector space over $\R$ with a positive definite scalar product $\langle \cdot, \cdot \rangle$, or over $\C$ with a positive definite hermitian product $\langle \cdot, \cdot \rangle$. Suppose $A$ is an operator on $V$. Let $\{v_1, v_2, \ldots, v_n\}$ be an orthonormal basis for $V$.
        \begin{enumerate}[label=(\alph*)]
            \item If $A$ is real-/complex-unitary, then $\{A v_1, A v_2, \ldots, A v_n\}$ is also an orthonormal basis for $V$.
            \item If $\{A v_1, A v_2, \ldots, A v_n\}$ is also an orthonormal basis for $V$, then $A$ is real-/complex-unitary.
        \end{enumerate}
    \end{propositionenv}
    \begin{proofenv}
        For (2), we can prove that $A$ is unitary by showing that $\| Av \| = \|v\|$ for any $v\in V$. Since $ \{v_1, v_2, \ldots, v_n\}$ is an orthonormal basis for $V$, we can write $v = \displaystyle \sum_{i=1}^n \alpha_i v_i$ for some unique $\alpha_i$ for $i=1,2,\ldots,n$. Then we have
        \[
        \| v \|^2 = \langle v, v \rangle = \sum_{i=1}^n |\alpha_i|^2, \quad \| Av \|^2 = \langle Av, Av \rangle = \sum_{i=1}^n |\alpha_i|^2.
        \]
        Thus $\| Av \| = \|v\|$ and $A$ is unitary.
    \end{proofenv}
    \begin{corollaryenv}
        Suppose $V$, $\C$, $\langle \cdot, \cdot \rangle$ satisfy the condition in Definition~\ref{def: complex unitary operator}. Then the following conditions on an operator $A$ on $V$ are equivalent:
        \begin{enumerate}
            \item $AA^* = A^*A$.
            \item $\left\| Av\right\| = \left\| A^* v\right\|$ for all $v\in V$.
            \item We can write $A=S+iT$ where $S$ and $T$ are self-adjoint operators on $V$ such that $ST=TS$. Moreover,
            $$ S = \frac{A+A^*}{2}, \quad T = \frac{A-A^*}{2i}. $$
        \end{enumerate}
    \end{corollaryenv}

    \newpage
